\chapter{Introduction}\label{introduction}
\chapterstatus{complete}

\section{What this project is}\label{introduction:section-aim}

The Math Project is a cumulative reference whose goal is to develop
mathematics, all of it, with complete proofs and in a single strict logical
order: the foundations (logic, set theory, the number systems, categories),
algebra, discrete mathematics, topology, analysis, probability and statistics,
differential equations and dynamical systems, numerical analysis, optimisation
and game theory, differential, complex and algebraic geometry, number theory,
and research-level subjects such as Hodge theory. Every subject is meant to be
treated in depth, in its natural generality and with the standard results that
a working mathematician expects to find, and at the same time slowly: each
definition is motivated and followed by examples, each hypothesis of a theorem
is shown to be necessary by a counterexample, proofs are divided into named
steps, and every chapter ends with exercises, most of them with solutions.

The organisation is modelled on the Stacks project: the text is written in a
strict logical order, every result has a permanent tag, and the whole project
is a single document that can be read as a book or browsed as a website. Like
the Stacks project, it is meant to be used by following references backwards:
any result can be traced, through the results its proof uses, down to the
axioms of set theory.

\begin{remark}\label{introduction:remark-history}
The project began in November 2026 from the chapters of
\href{https://lufranovittorio.github.io/Appunti_Matematica/}{The Hodge Project}
(source on GitHub: \url{https://github.com/lufranovittorio/Appunti_Matematica}), a
reference whose goal was to develop, with complete proofs, all the mathematics
needed to state and understand the Hodge conjecture. Informally, the Hodge
conjecture says that on a smooth complex projective variety $X$ every rational
cohomology class of type $(p,p)$ is a rational linear combination of the
classes of algebraic subvarieties of $X$ (Chapter~\ref{hodge-conjecture}).
Each word of this statement hides a substantial theory: rational cohomology
classes come from algebraic topology (Chapter~\ref{singular-homology}), the
type $(p,p)$ comes from the Hodge decomposition of a compact K\"ahler manifold
(Chapter~\ref{kahler}), which rests on the analysis of elliptic
operators (Chapter~\ref{elliptic-operators}), and the classes of subvarieties
come from algebraic geometry (Chapters~\ref{varieties} and~\ref{intersection-theory}).
This is why the parts on topology, analysis and geometry are, at present, the
most developed, and why the route to the Hodge conjecture is the most detailed
of the reading paths of Section~\ref{introduction:section-path}. The other
parts are being written in the order of the plan of the project; the chapters
that have not been written yet are present as outlines, with the status
\emph{planned} (Remark~\ref{introduction:remark-status}).
\end{remark}

\section{How the project is organised}\label{introduction:section-organisation}

The book has three levels. A \emph{part} is a large area of mathematics
(algebra, topology, number theory, \ldots); a \emph{chapter} is a region of
that area, a subject one could teach as a course (groups, homotopy theory,
analytic number theory, \ldots); a \emph{section} is a single topic inside it.
Chapters marked \emph{planned} consist, for now, of an outline of their
sections. The fifteen parts, grouped by the way they depend on each other, are
the following.

\medskip\noindent\textbf{Foundations and algebra.}
\begin{itemize}
\item Part~\ref{part-foundations}, \emph{Foundations}: logic, set theory, the
number systems, model theory, category theory, computability and complexity.
\item Part~\ref{part-algebra}, \emph{Algebra}: groups, rings, fields and Galois
theory, elementary number theory, linear and multilinear algebra,
representations of groups and Lie algebras, commutative algebra, homological
algebra and derived categories.
\item Part~\ref{part-discrete-mathematics}, \emph{Discrete Mathematics}:
combinatorics, generating functions, graph theory and the probabilistic
method.
\end{itemize}

\noindent\textbf{Topology.}
\begin{itemize}
\item Part~\ref{part-topology}, \emph{Topology}: general topology, homotopy and
covering spaces, topological manifolds, singular homology and cohomology,
Poincar\'e duality, cohomology operations, low-dimensional topology (surfaces,
knots, three-manifolds), sheaves and sheaf cohomology.
\end{itemize}

\noindent\textbf{Analysis and probability.}
\begin{itemize}
\item Part~\ref{part-measure-theory}, \emph{Real Analysis and Measure Theory}:
real analysis, measure and integration, differentiation and absolute
continuity, signed and complex measures with the duality of $L^p$ spaces, Radon
measures, and geometric measure theory (Hausdorff measures,
fractals, the area and coarea formulas, rectifiability).
\item Part~\ref{part-analysis}, \emph{Complex and Functional Analysis}: complex
analysis, special functions, functional analysis, several complex variables,
distributions, harmonic analysis and operator algebras.
\item Part~\ref{part-probability}, \emph{Probability and Statistics}:
probability theory, martingales and stochastic processes, statistics and
information theory.
\end{itemize}

\noindent\textbf{Number theory.}
\begin{itemize}
\item Part~\ref{part-number-theory}, \emph{Number Theory}, continuing the
elementary number theory of Chapter~\ref{number-theory}: prime numbers,
continued fractions and Diophantine approximation, the geometry of numbers,
analytic number theory (the prime number theorem, Dirichlet's theorem),
additive number theory and sieves, $p$-adic numbers and the local-global
principle, Diophantine equations, and then algebraic number theory: number
fields, local fields and class field theory.
\end{itemize}

\noindent\textbf{Equations and computation.}
\begin{itemize}
\item Part~\ref{part-differential-equations}, \emph{Differential Equations and
Dynamical Systems}: ordinary differential equations and dynamical systems,
ergodic theory, partial differential equations, the calculus of variations.
\item Part~\ref{part-numerical-analysis}, \emph{Numerical Analysis and
Optimisation}: numerical analysis, convex analysis and optimisation, game
theory.
\end{itemize}

\noindent\textbf{Geometry.}
\begin{itemize}
\item Part~\ref{part-differential-geometry}, \emph{Differential Geometry}: from
smooth manifolds and Lie groups through bundles, forms, connections and
characteristic classes to topological $K$-theory, Riemannian geometry,
currents and varifolds, symplectic geometry, elliptic operators, the Hodge
theorem for harmonic forms, spin geometry and the index theorem.
\item Part~\ref{part-complex-geometry}, \emph{Complex Geometry}: complex and
K\"ahler manifolds, the Hodge decomposition, the Lefschetz theorems, the
Kodaira vanishing and embedding theorems, deformations of complex structures and
moduli, and Calabi--Yau manifolds.
\item Part~\ref{part-algebraic-geometry}, \emph{Algebraic and Arithmetic
Geometry}: varieties, schemes, coherent cohomology, GAGA, curves, algebraic
groups and abelian varieties, algebraic cycles, algebraic surfaces, \'etale
cohomology, modular forms and the arithmetic of elliptic curves.
\item Part~\ref{part-hodge-theory}, \emph{Hodge Theory}: Hodge structures and
Mumford--Tate groups, the Hodge conjecture, the evidence for it, its failed
variants, mixed Hodge structures, absolute Hodge classes, and the Tate and
standard conjectures around it.
\end{itemize}

\noindent\textbf{Physics.}
\begin{itemize}
\item Part~\ref{part-mathematical-physics}, \emph{Mathematical Physics}:
classical mechanics, quantum mechanics, special and general relativity, as
applications of the analysis and geometry developed before.
\end{itemize}
The order of the chapters is the logical order of the project, not a
recommended reading order: a part comes after the parts it uses, so discrete
mathematics, which needs only algebra, comes before topology; probability,
which is built on measure theory, comes right after analysis; number theory,
which uses algebra and analysis but no geometry, comes next, from its
elementary and analytic chapters to algebraic number theory; the arithmetic of
curves (modular forms, elliptic curves) comes after algebraic geometry, whose
language it needs; and mathematical physics, which uses all of analysis and
geometry, comes last.
Section~\ref{introduction:section-path} describes how to read the book along
several reading paths.

\section{How to read this book}\label{introduction:section-path}

There are three ways to use the book.
\begin{itemize}
\item \emph{As a reference.} Look up a result by its tag or with the search page of the website, and follow its references backwards. The \textbf{Prerequisites} line at the start of
each chapter lists the chapters it depends on, and the dependency graph of the website shows, for every result, the results that its statement and proof use.
\item \emph{Along a reading path}, reading only what is needed to reach a chosen subject. The paths are described in
Subsection~\ref{introduction:subsection-reading-paths}; the longest and most detailed one leads to the Hodge conjecture
(Subsection~\ref{introduction:subsection-direct-route}).
\item \emph{As a course}, from the foundations onwards, in the order of the chapters. The book is written so that this is possible: nothing is used before it is proved
(Remark~\ref{introduction:remark-no-forward-references}).
\end{itemize}

\subsection{Reading paths}\label{introduction:subsection-reading-paths}

A reading path starts from the foundations and ends at a subject. Every path begins in the same way: Chapter~\ref{sets} for the language of sets, functions, cardinals and Zorn's
lemma, and Chapter~\ref{number-systems} for the integers, the rationals and the real and complex numbers (a reader who knows them can skim it). Chapter~\ref{logic} is needed only on
the paths that lead to logic itself, to model theory and to computability. The paths below list, after this common start, the chapters they use; inside a chapter, its introduction
says which sections are needed for what. Paths whose last part is still planned are described with the chapters already written that they will use.
\begin{itemize}
\item \emph{Algebra.} Chapters~\ref{groups}, \ref{rings}, \ref{fields} and~\ref{linear-algebra} are the common core; then multilinear algebra, representation theory and
commutative algebra in any order, and homological algebra, spectral sequences and derived categories in this order. Chapter~\ref{categories} is the language of the last three.
\item \emph{Elementary and algebraic number theory.} Chapters~\ref{groups}, \ref{rings} and~\ref{number-theory} (congruences, primitive roots, quadratic reciprocity, sums of
squares, quadratic forms and the Heegner numbers); then Part~\ref{part-number-theory}. In it, Chapter~\ref{prime-numbers} needs only real analysis besides these;
Chapters~\ref{continued-fractions} and~\ref{geometry-of-numbers} add measure theory (Chapter~\ref{measure-theory}) for their metric results and volumes;
Chapter~\ref{analytic-number-theory} needs complex analysis (Chapter~\ref{complex-analysis}), Poisson summation (Chapter~\ref{distributions}) and Fubini's theorem, and develops the
Gamma function itself; Chapter~\ref{p-adic-numbers} uses the previous two, and
Chapter~\ref{additive-number-theory} adds harmonic analysis (Chapter~\ref{harmonic-analysis}) for the circle method, and
Chapter~\ref{diophantine-equations} uses the elementary and analytic chapters. Galois theory (Chapter~\ref{fields}) and commutative algebra (Chapter~\ref{commutative-algebra}) then lead to
the algebraic chapters of the part (Chapters~\ref{algebraic-number-theory} and~\ref{local-fields}), and algebraic curves (Chapter~\ref{curves}) to the modular forms and elliptic curves of
Part~\ref{part-algebraic-geometry}.
\item \emph{Discrete mathematics.} Chapters~\ref{groups} (actions, for counting orbits), \ref{rings} (polynomial rings, which formal power series generalise)
and~\ref{linear-algebra} (for spectral graph theory), then Part~\ref{part-discrete-mathematics}. The probabilistic method uses only finite probability spaces, defined there.
\item \emph{Topology and geometry.} Chapters~\ref{general-topology}, \ref{homotopy} and~\ref{singular-homology}, then manifolds
(Chapters~\ref{topological-manifolds} and~\ref{smooth-manifolds}) and the rest of Part~\ref{part-differential-geometry}. Real analysis up to the inverse function theorem
(Chapter~\ref{real-analysis}) is needed from smooth manifolds onwards.
\item \emph{Low-dimensional topology.} Chapters~\ref{groups}, \ref{general-topology}, \ref{homotopy} (the fundamental group, van Kampen's theorem and covering spaces),
\ref{singular-homology} and~\ref{poincare-duality}, then Chapter~\ref{low-dimensional-topology} (surfaces and mapping class groups, knots and their polynomials, braids,
three-manifolds); the survey of four-manifolds there mentions intersection forms and gauge theory informally.
\item \emph{Analysis.} Chapters~\ref{real-analysis}, \ref{general-topology} (metric spaces, compactness and completeness), \ref{measure-theory}, \ref{lebesgue-differentiation},
\ref{signed-measures}, \ref{radon-measures}, \ref{complex-analysis}, \ref{functional-analysis} and~\ref{distributions}, with Chapter~\ref{linear-algebra} for spectral theory in finite dimension.
Chapter~\ref{special-functions} (Gamma, Bessel and hypergeometric functions, orthogonal polynomials, elliptic and theta functions, asymptotic methods) needs only real and complex
analysis and can be read right after Chapter~\ref{complex-analysis}.
\item \emph{Geometric measure theory.} Chapters~\ref{measure-theory}, \ref{lebesgue-differentiation} and~\ref{radon-measures}, then Chapter~\ref{geometric-measure-theory}
(Brunn--Minkowski and isoperimetric inequalities, Hausdorff measures and dimension, self-similar fractals, Rademacher's theorem, the area and coarea formulas, densities and rectifiable
sets, Frostman's lemma and projections); its results on Lipschitz maps and rectifiable sets are used again on manifolds (Chapter~\ref{smooth-manifolds}). The path continues, after
Chapters~\ref{multilinear-algebra} (exterior algebra and comass) and~\ref{de-rham} (forms, Stokes' theorem, de Rham currents), with Chapter~\ref{currents} (mass, normal and integral currents, Plateau's
problem, calibrations; varifolds, first variation and monotonicity).
\item \emph{Probability and statistics.} The analysis path up to Chapter~\ref{radon-measures} (measures, the integral, $L^p$ spaces, modes of convergence, independence and the
Borel--Cantelli lemmas, Fubini's theorem in Chapter~\ref{measure-theory}, the Radon--Nikodym theorem in Chapter~\ref{lebesgue-differentiation}, uniform integrability in Chapter~\ref{signed-measures}, weak convergence of measures in
Chapter~\ref{radon-measures}), Hilbert spaces from Chapter~\ref{functional-analysis} (conditional
expectation is an orthogonal projection), and the Fourier transform of Chapter~\ref{distributions} (characteristic functions); then Part~\ref{part-probability}.
\item \emph{Differential equations and dynamics.} The analysis path, together with the existence theorem of Picard and Lindel\"of in Chapter~\ref{real-analysis} and the Jordan
normal form of Chapter~\ref{linear-algebra}; then Part~\ref{part-differential-equations}. Flows on manifolds use Chapter~\ref{smooth-manifolds}, read informally until
Part~\ref{part-differential-geometry}.
\item \emph{Numerical analysis.} Chapters~\ref{real-analysis} and~\ref{linear-algebra}, then Part~\ref{part-numerical-analysis}; its second half uses
Part~\ref{part-differential-equations}.
\item \emph{Optimisation and game theory.} Chapter~\ref{real-analysis} (whose last sections already contain Lagrange multipliers, Farkas' lemma and the Karush--Kuhn--Tucker
conditions), compactness from Chapter~\ref{general-topology}, convexity and separation from Chapter~\ref{functional-analysis}, and the Brouwer fixed point theorem from
Chapter~\ref{singular-homology} (for Nash equilibria); then Part~\ref{part-numerical-analysis}.
\item \emph{Mathematical physics.} For classical mechanics: Chapters~\ref{real-analysis}, \ref{linear-algebra}, \ref{ordinary-differential-equations}
and~\ref{calculus-of-variations}, then manifolds, Lie groups and Riemannian and symplectic geometry from Part~\ref{part-differential-geometry}
(Chapters~\ref{smooth-manifolds}, \ref{lie-groups}, \ref{riemannian-geometry} and~\ref{symplectic-geometry}). For quantum mechanics: the analysis path with
Chapters~\ref{harmonic-analysis} and~\ref{operator-algebras} (spectral theory of unbounded operators), Chapter~\ref{representation-theory} (the representations of
$\mathrm{SU}(2)$) and Chapter~\ref{probability}. For relativity: Chapters~\ref{connections} and~\ref{riemannian-geometry}, with hyperbolic equations from
Chapter~\ref{sobolev-spaces}. Then Part~\ref{part-mathematical-physics}.
\item \emph{The Hodge conjecture.} The path of the original
\href{https://lufranovittorio.github.io/Appunti_Matematica/}{Hodge Project}, described in detail in Subsection~\ref{introduction:subsection-direct-route}.
\end{itemize}

\subsection{The route to the Hodge conjecture}\label{introduction:subsection-direct-route}

The route to the Hodge conjecture reads only what is needed to understand its statement and the one case that is proved in full, the Lefschetz $(1,1)$ theorem: it ends with
Chapter~\ref{hodge-conjecture}, which states the conjecture and, in Section~\ref{hodge-conjecture:section-lefschetz-11}, proves it for divisors. Within a chapter, the sections that the route
actually uses can be found by following references backwards from these two chapters; on the website, the dependency graph of a result shows them directly. A reader who
skips a section and later meets a reference to it can go back. Part by part, the route is as follows.
\begin{itemize}
\item \emph{Foundations.} Read Chapter~\ref{sets} for the language of sets, functions, cardinals and Zorn's lemma, and Chapter~\ref{categories} for the language of categories,
functors, limits and adjoint functors. A reader who knows the real and complex numbers can skim Chapter~\ref{number-systems}. Chapters~\ref{logic} and~\ref{model-theory} are
not needed.
\item \emph{Algebra.} Linear and multilinear algebra are needed in full: the complex structures of Chapter~\ref{linear-algebra} and the exterior algebra and forms of type $(p,q)$ of
Chapter~\ref{multilinear-algebra} are the linear algebra of the Hodge decomposition. From Chapters~\ref{rings}, \ref{fields} and~\ref{commutative-algebra} one needs ideals,
modules and tensor products, transcendence degree, and the commutative algebra behind dimension and smoothness of varieties. Homological algebra and the spectral sequences of filtered
and double complexes are needed for sheaf cohomology; derived categories only for hypercohomology. From Chapter~\ref{representation-theory} only the representations of
$\mathfrak{sl}_2$ are used, for hard Lefschetz. Chapter~\ref{number-theory}, most of Chapter~\ref{groups}, Galois theory and the structure of semisimple Lie algebras can be
postponed.
\item \emph{Topology.} General topology up to partitions of unity; homotopy, the fundamental group and covering spaces; singular homology and cohomology with cup and cap products;
Poincar\'e duality and fundamental classes. Chapter~\ref{cohomology-operations} is not needed.
\item \emph{Sheaves.} Both chapters: sheaves and their images, sheaf cohomology through flasque and soft resolutions, \v{C}ech cohomology and the comparison with singular cohomology.
\item \emph{Analysis.} Real analysis up to the inverse and implicit function theorems; the Lebesgue integral, $L^p$ spaces and Fubini's theorem; Cauchy's theorem in one and several
complex variables, and the Dolbeault--Grothendieck lemma; Banach and Hilbert spaces and compact operators; Sobolev spaces. The geometric measure theory at the end of
Chapter~\ref{measure-theory} and the second half of Chapter~\ref{functional-analysis} are not on the route.
\item \emph{Differential geometry.} Smooth manifolds, vector bundles, differential forms and Stokes' theorem, de Rham's theorem, connections and curvature, Chern--Weil theory,
the Hodge star, elliptic operators (whose estimates can be taken on faith on a first reading) and the Hodge theorem. Chapter~\ref{lie-groups}, most of
Chapter~\ref{riemannian-geometry} and Chapter~\ref{symplectic-geometry} can be postponed.
\item \emph{Complex geometry.} Complex manifolds, holomorphic bundles, Dolbeault cohomology, K\"ahler manifolds, the Hodge decomposition, hard Lefschetz and the Hodge--Riemann
relations, positive line bundles and the Kodaira theorems. Chapters~\ref{deformations} and~\ref{calabi-yau} are not needed.
\item \emph{Algebraic geometry.} Varieties; schemes, mainly as a language; coherent sheaves and their cohomology; divisors and line bundles; GAGA, which makes smooth projective
varieties into compact K\"ahler manifolds; algebraic cycles and the cycle class map. Chapters~\ref{curves}, \ref{algebraic-groups}, \ref{intersection-theory}
and~\ref{etale-cohomology} are not needed; from Chapter~\ref{algebraic-groups} only the Deligne torus is used.
\item \emph{Hodge theory.} The definitions of Chapter~\ref{hodge-structures}, then Chapter~\ref{hodge-conjecture}. The remaining chapters of the part are the
surrounding theory, to be read afterwards.
\end{itemize}

\subsection{The main line}\label{introduction:subsection-main-line}

Seen from above, the direct route has the following stages; everything else in the book supports them or grows out of them.
\begin{enumerate}
\item \emph{Cohomology.} Singular cohomology and Poincar\'e duality
(Chapters~\ref{singular-homology} and~\ref{poincare-duality}), sheaf
cohomology (Chapter~\ref{sheaves}) and de Rham's theorem
(Chapter~\ref{de-rham}) give the rational cohomology of a manifold and its
description by differential forms.
\item \emph{Harmonic theory.} Elliptic operators
(Chapter~\ref{elliptic-operators}) give the Hodge theorem: every de Rham class
has a unique harmonic representative (Chapter~\ref{hodge-theorem}).
\item \emph{K\"ahler geometry.} On a compact K\"ahler manifold
(Chapter~\ref{kahler}) the Laplacian respects the bidegree of forms, which gives
the Hodge decomposition $H^k = \bigoplus_{p+q=k}H^{p,q}$
(Section~\ref{kahler:section-hodge-decomposition}), hard Lefschetz and the Hodge--Riemann
relations (Section~\ref{kahler:section-lefschetz-theorems}).
\item \emph{Algebraic geometry.} Projective varieties are K\"ahler, and GAGA
identifies analytic and algebraic objects on them (Chapter~\ref{gaga}).
Subvarieties have cohomology classes, which are rational and of type $(p,p)$
(Chapter~\ref{intersection-theory}).
\item \emph{The conjecture.} Hodge structures (Chapter~\ref{hodge-structures})
isolate the linear algebra, and the Hodge conjecture asserts the converse of
the last statement (Chapter~\ref{hodge-conjecture}). It is a theorem for
divisors (Chapter~\ref{hodge-conjecture}), it is known in further cases
(Chapter~\ref{known-cases}), and its natural strengthenings are false
(Chapter~\ref{known-cases}).
\item \emph{The surrounding theory.} Variations of Hodge structure and normal
functions (Chapter~\ref{vhs}), mixed Hodge
structures (Chapter~\ref{mixed-hodge-structures}), absolute Hodge classes
(Chapter~\ref{absolute-hodge}), and the Tate and standard conjectures
(Chapter~\ref{absolute-hodge}) place the
conjecture in its wider landscape.
\end{enumerate}

\begin{remark}\label{introduction:remark-first-reading}
Most of the book is not needed on the way to the Hodge conjecture: everything in the parts on discrete mathematics, probability, differential equations, numerical analysis,
optimisation and number theory, and, in the older parts, model theory and elementary number theory, geometric measure theory, Riemannian and symplectic geometry for their own
sake, Calabi--Yau manifolds, curves and abelian varieties. Each of these can be read at any time after its prerequisites.
\end{remark}

\section{A map of the integrals}\label{introduction:section-integration}

Many integrals appear in the book: the Riemann integral, the Lebesgue integral, volumes, lengths and areas, line integrals of two kinds, complex line integrals, integrals of differential forms,
integrals on Riemannian manifolds and on Lie groups. They are defined where they are first needed, often by an explicit formula, but they are not independent notions. All of them are special cases
of two constructions, and the second is built from the first.
\begin{enumerate}
\item \emph{The integral of a function against a measure}, $\int_Xf\,d\mu$ (Chapter~\ref{measure-theory}, Section~\ref{measure-theory:section-integral}). It needs a measure on $X$, that is, a way to assign a size to
subsets, and no orientation: reversing a curve does not change its length. Integrals of this kind are written with $d\mu$, $dx$, $ds$, $dA$ or $d\mathcal{H}^k$.
\item \emph{The integral of a differential $k$-form over an oriented $k$-dimensional manifold}, $\int_M\omega$ (Definition~\ref{differential-forms:definition-integral}). It needs an orientation, and no measure or metric:
in a positively oriented chart, $\omega = f\,dx^1 \wedge \cdots \wedge dx^k$ and $\int_M\omega$ is the Lebesgue integral $\int f\,d\lambda^k$, patched together with a partition of unity. Reversing the orientation
changes the sign.
\end{enumerate}
On an oriented Riemannian manifold the two meet: a $k$-form is a function times the volume form, and $\int_M\omega = \int_M\langle\omega, \xi\rangle\,d\mu_g$ with the orienting $k$-vector $\xi$
(Proposition~\ref{riemannian-geometry:proposition-orienting-vector}). Here is how the individual integrals of the book fit into this picture.
\begin{itemize}
\item \emph{The Riemann integral} $\int_a^bf$ (Definition~\ref{real-analysis:definition-riemann}) comes first because it needs no measure theory, and it is enough for continuous functions on intervals. It is the
Lebesgue integral $\int_{[a,b]}f\,d\lambda$ for every Riemann integrable $f$ (Propositions~\ref{measure-theory:proposition-riemann-lebesgue} and~\ref{measure-theory:proposition-riemann-integrable}), so nothing proved
with it has to be redone. Vector- and complex-valued functions are integrated componentwise, $\int(u + iv) = \int u + i\int v$.
\item \emph{Multiple integrals and volumes} $\int_{\RR^n}f\,d\lambda^n$: Lebesgue measure on $\RR^n$, computed by Fubini's theorem and the change of variables formula
(Sections~\ref{measure-theory:section-fubini} and~\ref{measure-theory:section-change-of-variables}), with Cavalieri's principle and solids of revolution as special cases
(Subsection~\ref{measure-theory:subsection-volumes}).
\item \emph{Length and line integrals of the first kind} $\int_\gamma f\,ds$ (Section~\ref{real-analysis:section-curves}): defined by the formula $\int_a^bf(\gamma(t))|\gamma^{\prime}(t)|\,dt$. They are integrals of the first
construction: for an injective regular curve, $\int_\gamma f\,ds = \int_{\gamma([a,b])}f\,d\mathcal{H}^1$, the integral against one-dimensional Hausdorff measure on the image (Theorem~\ref{geometric-measure-theory:theorem-area-formula} and
Example~\ref{geometric-measure-theory:example-area-formula}(1)), which depends only on the set.
\item \emph{Areas and surface integrals} $\int_\Sigma f\,dA$: the integral against $\mathcal{H}^2$, or $\mathcal{H}^k$ for $k$-dimensional submanifolds, computed from a parametrisation by the area formula
(Section~\ref{geometric-measure-theory:section-area-formula}); equivalently, the Riemannian measure of the induced metric (Propositions~\ref{riemannian-geometry:proposition-riemannian-measure}
and~\ref{riemannian-geometry:proposition-hausdorff-measure}). The theorems of Pappus--Guldin are examples.
\item \emph{Line integrals of the second kind} $\int_\gamma F \cdot d\gamma$ (Definition~\ref{real-analysis:definition-second-kind}): the integral of the $1$-form $F_1\,dx_1 + \cdots + F_n\,dx_n$ over the oriented
curve (Definition~\ref{differential-forms:definition-line-integral}), an integral of the second construction. For a regular curve it is $\int_\gamma\langle F, T\rangle\,ds$, a first-kind integral of the tangential component
(Proposition~\ref{real-analysis:proposition-second-kind}(3)): choosing the unit tangent $T$ is choosing an orientation.
\item \emph{Complex line integrals} $\int_\gamma f\,dz$ (Definition~\ref{complex-analysis:definition-path-integral}): the integral of the complex-valued $1$-form $f\,dz$, which splits into two real line integrals of the
second kind, $\int_\gamma(u\,dx - v\,dy) + i\int_\gamma(v\,dx + u\,dy)$ (Proposition~\ref{complex-analysis:proposition-path-integral-real}).
\item \emph{Flux integrals and integrals of $k$-forms} over oriented submanifolds, $\int_M\omega$, with Stokes' theorem (Theorem~\ref{differential-forms:theorem-stokes}) and its classical cases, the theorems of Green, Gauss and
Kelvin--Stokes (Chapter~\ref{riemannian-geometry}): the second construction.
\item \emph{Integrals on Riemannian manifolds} $\int_Mf\,d\mu_g$, and on Lie groups against Haar measure (Theorem~\ref{lie-groups:theorem-haar}): the first construction, for the Riemannian measure and the Haar measure. On an
oriented manifold, $\int_Mf\,d\mu_g = \int_Mf\,\mathrm{vol}_g$.
\item \emph{Integration in cohomology}: $\int_M \colon H^n_{\mathrm{dR}}(M) \to \RR$ on a compact oriented manifold, periods of closed forms over cycles, and the integrals of characteristic forms that compute Chern
numbers: the second construction applied to closed forms, which depends only on cohomology classes by Stokes' theorem (Chapters~\ref{de-rham} and~\ref{characteristic-classes}).
\item \emph{Distributions and currents} (Chapter~\ref{distributions} and Section~\ref{currents:section-currents-manifold}) are not integrals but generalise them: a function $f$ gives the functional $\varphi \mapsto \int f\varphi$
and a submanifold $Z$ the functional $\varphi \mapsto \int_Z\varphi$, and the theory keeps only the functional.
\end{itemize}
In short: if an integral ignores orientation, look for the measure; if it changes sign when the orientation is reversed, it is the integral of a form. The Lebesgue integral underlies both.

\section{Sources and status}\label{introduction:section-sources}

The mathematics in this project is standard, and each chapter lists the
references it follows; the bibliography collects them. The chapters of
Part~\ref{part-hodge-theory} follow the original papers and the surveys cited in
each of them. Where a theorem is quoted
rather than proved, the text says so and gives a reference
(Remark~\ref{introduction:remark-proofs}). The chapters inherited from
\href{https://lufranovittorio.github.io/Appunti_Matematica/}{The Hodge Project} are written, and those of the new parts are being written; the
status of every chapter is shown at its start on the website and in the table
of contents. No chapter has yet been checked by an independent reader, which is
the condition for the status \emph{checked}
(Remark~\ref{introduction:remark-status}).

\begin{remark}\label{introduction:remark-depth}
The standard of depth is that of the Stacks project: a chapter on a subject
should contain the definitions and results that a specialist expects to find
there, in their natural generality, with complete proofs, and not only a
selection. A chapter is \emph{complete} when it meets this standard. Around
the core, chapters also collect the well-known results that a curious reader
looks for (in number theory, for example, the Gaussian integers, the Heegner
numbers and Euler's polynomial $n^2 + n + 41$), even when nothing later uses
them.
\end{remark}

\section{Conventions of the text}\label{introduction:section-conventions}

\begin{notation}\label{introduction:notation-numbering}
All numbered items of a section share a single counter: Lemma~5.2.3 is the
third numbered item of Section~5.2, and it is followed by Example~5.2.4.
Chapters are numbered consecutively through all the parts.
\end{notation}

\begin{notation}\label{introduction:notation-environments}
The text uses the following kinds of numbered items.
\begin{itemize}
\item \textbf{Definition}: introduces a notion.
\item \textbf{Theorem}, \textbf{Proposition}, \textbf{Lemma},
\textbf{Corollary}: proved statements, in roughly decreasing order of
importance.
\item \textbf{Conjecture}: a statement that is not known to be true.
\item \textbf{Example}: an instance of a definition or a result.
\item \textbf{Counterexample}: shows that a hypothesis cannot be dropped, or
that a plausible statement is false.
\item \textbf{Remark}: context, motivation or intuition. Remarks are never used
in proofs.
\item \textbf{Notation}: fixes notation or a convention.
\item \textbf{Exercise}: a statement for the reader to prove, sometimes followed
by a solution.
\end{itemize}
\end{notation}

\begin{remark}\label{introduction:remark-proofs}
A result followed by \emph{Proof} has a complete proof that uses only earlier
material. A result followed by \emph{Sketch of proof} has an argument whose
main ideas are given but whose omitted verifications are named. A result
followed by no proof is quoted from the literature, with a reference.
Sketches and quotations are temporary: the aim of the project is to replace
all of them with complete proofs.
\end{remark}

\begin{remark}\label{introduction:remark-no-forward-references}
The project has no forward references: a proof may only use definitions and
results that appear before it. Examples, counterexamples, remarks and
exercises may mention notions that are defined later (for instance, the
category of topological spaces appears as an example in
Chapter~\ref{categories}), but such mentions are never used in a proof. The
only exception is this introduction, which refers forward freely.
\end{remark}

\begin{remark}\label{introduction:remark-tags}
Every labelled item (chapter, section, definition, result, example, and so on)
carries a permanent four-character \emph{tag}, shown next to it on the
website. The number of an item changes whenever material is inserted before
it; its tag never changes. To refer to an item from outside the project, cite
its tag.
\end{remark}

\begin{remark}\label{introduction:remark-status}
Each chapter has a status. \emph{Planned}: only an outline exists.
\emph{Draft}: the chapter is being written, and its proofs have not been
checked. \emph{Complete}: all planned material is written, and the proofs
await independent checking. \emph{Checked}: every proof in the chapter has been
checked by a reader other than its author.
\end{remark}

\section{Mathematical conventions}\label{introduction:section-mathematical-conventions}

\begin{notation}\label{introduction:notation-universe}
We work in ZFC set theory together with the assumption that a Grothendieck
universe exists, and we fix such a universe $\mathcal{U}$ once and for all. A
set is \emph{small} if it is an element of $\mathcal{U}$. Chapter~\ref{sets}
develops these notions; Chapter~\ref{categories} explains why they are needed.
\end{notation}

\begin{notation}\label{introduction:notation-numbers}
The natural numbers include zero: $\NN = \{0, 1, 2, \ldots\}$. This is the
convention of set theory and logic, where $\NN$ is built starting from
$0 = \emptyset$ and each natural number $n = \{0, 1, \ldots, n - 1\}$ is a set with
$n$ elements (Remark~\ref{sets:remark-zero-natural}). Many texts, especially in
analysis and number theory, start from $1$ instead; for the positive integers we
write $\NN_{>0} = \{1, 2, 3, \ldots\}$, and more generally
$\NN_{\geq k} = \{k, k + 1, \ldots\}$. A real number $x$ is \emph{positive} if
$x > 0$ and \emph{non-negative} if $x \geq 0$.
\end{notation}

\begin{notation}\label{introduction:notation-maps}
The composition of maps $f \colon X \to Y$ and $g \colon Y \to Z$ is written
$g \circ f$ or $gf$. The identity map of $X$ is written $\id_X$.
\end{notation}

\begin{notation}\label{introduction:notation-rings}
A ring is associative and has an identity element $1$, and ring homomorphisms
preserve identity elements. Rings are not assumed to be commutative, except
from Chapter~\ref{commutative-algebra} onwards, where ``ring'' means
``commutative ring'' unless stated otherwise.
\end{notation}
